%!TEX root = ../main.tex
%% ===============================================================
%% 기호·연산자 매크로 (내용 계층). 레이아웃은 pnuthesis.cls, 수치는 config/numbers.tex.
%% 용어 고정 (docs/STYLE.md): 교전 / 한타 T / 소규모 교전 S / 변화 방향 / 교전 전 상태 / 승률·시간 스플라인 기준선.
%% ===============================================================
\newcommand{\ind}{\mathbf{1}}
\newcommand{\R}{\mathbb{R}}
\newcommand{\E}{\mathbb{E}}
\DeclareMathOperator*{\argmin}{arg\,min}
\DeclareMathOperator*{\argmax}{arg\,max}
\DeclareMathOperator{\Brier}{Brier}
\DeclareMathOperator{\MCB}{MCB}
\DeclareMathOperator{\DSC}{DSC}
\DeclareMathOperator{\UNC}{UNC}
\DeclareMathOperator{\AUC}{AUC}
\newcommand{\todo}[1]{\textcolor{red}{[TODO: #1]}}
\newcommand{\pendingauthor}[1]{\textcolor{blue}{[저자 확인: #1]}}
%% 표주: 캡션은 이름만 쓰고, 가중·보정기·분모·구간 수준 같은 최소 조건은 표 아래 한 줄에 둔다 (docs/STYLE.md).
\newcommand{\tabnote}[1]{\par\vspace{3pt}\begin{minipage}{\linewidth}\footnotesize\setstretch{1.1}#1\end{minipage}}

%% --- 연구 대상
\newcommand{\Vhat}{\widehat V}                       % 동결 승률 평가기
\newcommand{\Vfold}[1]{\widehat V^{(-#1)}}          % fold 제외 평가기
\newcommand{\dV}{\Delta\widehat V}                   % 구간 변화
\newcommand{\SVI}{\mathrm{SVI}}
\newcommand{\Ysvi}{Y_{\mathrm{SVI}}}
\newcommand{\ppre}{p_{\mathrm{pre}}}
\newcommand{\Spre}{S_{\mathrm{pre}}}
\newcommand{\xpre}{x_{\mathrm{pre}}}
\newcommand{\xpost}{x_{\mathrm{post}}}
\newcommand{\tcut}{\tau}                             % 예측 시점 (첫 킬 − 15 s)
\newcommand{\eh}{e_{h}}                              % 결과 시점
\newcommand{\PTflex}{\mathrm{PT}_{\mathrm{flex}}}
\newcommand{\PTlin}{\mathrm{PT}_{\mathrm{linear}}}
\newcommand{\PTflexS}{\mathrm{PT}_{\mathrm{flex},S}}
\newcommand{\PTlinS}{\mathrm{PT}_{\mathrm{linear},S}}
\newcommand{\qS}{q_{S}}
\newcommand{\qTS}{q_{TS}}
\newcommand{\qTtoS}{q_{T\to S}}
\newcommand{\Bforty}{\mathrm{B40}}
\newcommand{\nmin}{n_{\min}}
\newcommand{\dBrier}{\Delta\mathrm{Brier}}
\newcommand{\ci}[2]{[#1,\ #2]}                       % 95% 구간 표기
